\section{Les interprétations réalistes}
\subsection{Brisure de la chaîne de Von Neumann (mécanisme dans les règles quantiques)}
\subsubsection{Décohérence macroscopique}
(Zurek, Zeh, Gell-Mann, Hartle, voir aussi la logique partielle de Griffiths, Omnès)
Considérons un atome dans un état quantique superposé du type
\[
a\ket{\alpha}+b\ket{\beta}
\]
tel que permis par l’Équation de E.Schrödinger. On parle alors d’état cohérent dans le sens
où les coefficients, a et b, sont reliés par la normalisation de l’état global et ont une phase
physique relative. Ces états ortho-normalisés $\ket{\alpha}$,
 $\ket{\beta}$ de l’atome peuvent être des états
propres d’énergie (excitation : modèle en couche), de position,... Supposons maintenant que
cet atome émet un photon dont l’impulsion correspond à un certain état de l’atome. L’état du
Système atome + photon pourrait alors s’écrire comme suit,
\[
a\ket{\alpha}c_{\alpha 1}\ket{k_\alpha}+b\ket{\beta}c_{\beta 1}\ket{k_\beta}
\]
$\ket{k_\alpha}$, $\ket{k_\beta}$ étant les états propres orthonormés de l’impulsion du photon
 et $c_{\alpha 1}$ , $c_{\beta 1}$ des coefficients ; il s’agit là d’un état dit intriqué. Si
 le photon a l’impulsion $k_\alpha$ ($k_\beta$), suite à une
mesure, alors l’atome est dans l’état réduit $\ket{\alpha}$ ($\ket{\beta}$) – ce qui fait
 sens physiquement ici.
Dans le cas où un deuxième photon serait émis par l’atome, il faudrait alors considérer l’état
global initial du Système atome + 2 photons – vis-à-vis de la mesure éventuelle des
impulsions des 2 photons émis – qui serait le suivant,
\[
a\ket{\alpha}c_{\alpha 1}\ket{k_\alpha}c_{\alpha 2}\ket{k_\alpha}+
b\ket{\beta}c_{\beta 1}\ket{k_\beta}c_{\beta 2}\ket{k_\beta}
\]
Une mesure macroscopique en physique classique revient justement à un grand nombre
(dans l’espace et le temps) d’interactions microscopiques et de mesures (quantiques) via par
exemple la lecture « à l’oeil » d’un cadran gradué muni d’une aiguille. Dans l’exemple
typique ci-dessus, où l’atome ferait partie du système physique soumis à la mesure
macroscopique, les nombreuses mesures s’appliqueraient aux photons « visibles » et
seraient associées à l’état initial du Système atome + N photons (avec N >> 1) :
\[
a\ket{\alpha}c_{\alpha 1}\ket{k_\alpha}c_{\alpha 2}\ket{k_\alpha}...c_{\alpha N}\ket{k_\alpha}+
b\ket{\beta}c_{\beta 1}\ket{k_\beta}c_{\beta 2}\ket{k_\beta}...c_{\beta N}\ket{k_\beta}
\]
Faisons l’hypothèse simplificatrice, tout en respectant la normalisation d’état, que par
exemple les |$c_{\alpha i}$| ~ 1 alors que les |$c_{\beta 1}$| << 1 [i = 1,...,N] (des valeurs
 autres conduiraient
simplement à une convergence de probabilité plus lente et possiblement vers l’autre état
final). Ainsi, l’amplitude de probabilité d’obtenir le résultat $k_\beta$ pour tous les
 photons tend vers zéro :
\[
A_\beta=c_{\beta 1}c_{\beta 2}c_{\beta 3}c_{\beta 4}c_{\beta 5}c_{\beta 6}... \rightarrow 0
\]
et on observe bien que dans cette limite (de système ‘augmenté’) les photons visibles
peuvent être dans l’état $\ket{\alpha}$, ce qui est intuitif (une certitude de signal
 d’information
croissante sur l’ensemble incluant les dernières émissions de photons) et signifie que
l’atome est lui dans l’état, maintenant décohérent,
associé à la valeur propre unique et classique, $\alpha$. L’on retrouve alors, dans cette
 limite du
formalisme quantique, le déterminisme de la mesure macroscopique (pour l’ensemble des
atomes du système macroscopique) en physique classique.
\subsubsection{Pragmatisme des laboratoires}
Une possibilité est qu’il existe un mécanisme, non encore inclus dans la version standard de
la Mécanique quantique (c.f. II.1), qui force l’annulation de l’amplitude ci-dessus,
\[
A_\beta=c_{\beta 1}c_{\beta 2}c_{\beta 3}c_{\beta 4}c_{\beta 5}...c_{\beta N} = 0
\]
au-delà d’un seuil de valeur de A$_\beta$ , suffisamment petit devant 1, donc après un certain
nombre N d’interactions de l’atome enrichissant le vecteur d’état initial. Cette annulation
réduit alors l’état de l’atome à
\[
\ket{\alpha}
\]
Cette version de la MQ évite donc la notion de superposition. L’existence d’interactions du
système pourrait en effet être inévitable, intrinsèque à la théorie, en tout cas de manière
effective et pragmatique au sein du laboratoire encadrant la mesure quantique. L’annulation
de A$_\beta$ met un terme à la régression infinie d’états de particules supplémentaires
 impliquées
dans les interactions successives, appelée chaîne de J.Von Neumann (introduite par ce
dernier afin de décrire les interactions du système avec l’appareil de mesure), au sens où les
nouveaux états s’ajoutant n’enrichissent plus un état intriqué. L’annulation de A$_\beta$ , ou
réduction de l’état de l’atome (existence d’une valeur déterminée pour une observable),
peut ({\footnotesize Si c’est la prise de conscience de l’observateur (possiblement via l’effet
 des courants électriques le long des neurones) qui brise la chaîne de Von Neumann, alors il
 s’agit de « l’interprétation de Wigner » (la réduction du
vecteur d’état est également liée à la conscience humaine et la faculté d’introspection chez
 London \& Bauer).}) avoir lieu en amont et donc indépendamment d’un observateur, de sorte que
  cette
approche peut se classer, d’un point de vue philosophique, dans le réalisme. Et la valeur
physique unique fait partie d’une vérité accessible (c.f. 1er tableau) puisque mesurable
[depuis le monde perçu et via notre faculté de pensée].
\subsection{Le réel voilé (sans modification des lois quantiques)}
Pour B.d’Espagnat, il existe un réel fondamental ne nécessitant pas d’être rapporté à
l’Homme, et, donc indépendant de notre appareil sensitif corporel (dont l’appareillage
expérimental est le prolongement). Cet auteur se situe ainsi au coeur du réalisme et se
rapproche en particulier de la vision du potentiel limité de la raison selon E.Kant en ce que,
pour le premier, seules certaines structures générales du réel sont au mieux accessibles à la
connaissance humaine. D’où l’appellation de réel voilé pour cette vision ontologique de la
Mécanique quantique. Un vecteur d’état superposé ne retranscrit ici au mieux qu’une vision
partielle du réel, évitant ainsi tout paradoxe à son égard au prix d’une forme d’ignorance
inéluctable.
\subsection{Les variables cachées (ajout d’un potentiel)}
Parmi les théories (voir par exemple Wiener \& Siegel, Hiley, Holland ou Landé) à ‘variables
cachées’ – dites aussi supplémentaires au sens non présentes dans la description standard de
la MQ (c.f. II.1) – celle de De Broglie-Bohm est sans doute la plus célèbre. Celle-ci confère
un rôle central à l’environnement d’un système physique. Il faut alors définir un état
quantique global du Système + Environnement, c’est-à-dire un état de type produit tensoriel
décrivant le système lui-même ainsi que chacun des nombreux éléments constituant son
environnement direct. Toutes ces composantes introduisent de nouvelles variables
quantiques les caractérisant que sont notamment leurs positions. Les vitesses associées (v),
dites vitesses quantiques, sont définies à partir du courant de probabilité (J) impliquant les
fonctions d’onde ($\Psi$) non normalisées (vitesses de phase), la normalisation se faisant sur
l’état global :
\[
J\wt{=}-|\Psi|^2\frac{dR}{dt}\wt{=}\rho v
\]
t dénotant le temps et $\rho$ la densité de probabilité de présence, dans le cas à 1 dimension
d’espace de coordonnée R. On retrouve ici la vitesse locale de la version hydrodynamique
de l’Équation de E.Schrödinger introduite par E.Madelung (en 1927). L’effet de
l’environnement (masse générique m) sur le Système intervient dans l’Équation de
Schrödinger via l’énergie potentielle additionnelle suivante (sommée sur toutes les
composantes),
\[
V=\frac{-\hbar^2}{2m}\frac{\Delta|\Psi(R,t)|}{|\Psi(R,t)|}
\]
où $\Delta$ est l’opérateur laplacien. Ce « potentiel quantique » permet de retrouver les
prédictions probabilistes de la description standard de la MQ. Ces probabilités ne font que
traduire notre niveau de méconnaissance de l’état du Système pré-existant à la procédure de
mesure, celui-ci étant fixé au sein de l’état global unique du Système + Environnement via
les interactions Système $\leftrightarrow$ Environnement (superposition effective et non
 intrinsèque). Le
résultat de la mesure lui ne fait que révéler, mettre en évidence la valeur propre unique
associée à cet état propre d’observable antérieur du Système : cette théorie quantique
recouvre le déterminisme classique. Ainsi, le système possède des propriétés bien
déterminées (position, impulsion, etc) indépendamment de l’observateur, ce qui peut placer
cette théorie dans le courant réaliste, avec accession à la vérité (c.f. 1er tableau) – par la
simple mesure expérimentale.
\subsection{Interprétation modale (ajout de vecteur d’état)}
L’Équation de Schrödinger guide l’évolution temporelle des états possiblement superposés
d’un système physique, appelés ici « états dynamiques ». Un autre état, dénommé ici « état
des valeurs » et ne pouvant être que pur par construction, pourrait lui correspondre à la
valeur propre unique mesurable ; l’état des valeurs représente les propriétés physiques du
système considéré. On ne peut atteindre cet état déjà attribué à un système qu’au moment de
la mesure de ses observables, ce qui classe possiblement cette interprétation dite modale
dans le réalisme, avec accessibilité à la vérité (c.f. 1er tableau) via une simple mesure en
laboratoire. Le formalisme modal permet, comme il se doit, de reconstruire les règles des
probabilités de Born afin de ne pas générer de conflit avec l’ensemble des jeux de données
expérimentales connus – probabilités reflétant ici juste notre ignorance partielle. On doit
cette solution astucieuse au problème de la superposition, ne lui conférant qu’un statut
formel d’implémentation temporelle, en particulier à B. van Fraassen ; d’autres versions de
l’interprétation modale (Kochen, Dieks et Healey) relient les états dynamiques et les états
des valeurs grâce à la décomposition bi-orthonormale (décomposition de Schmidt).
\subsection{Modification de l’Équation de Schrödinger (ajout d’un terme)}
En plus de décrire la dynamique temporelle, l’Équation de Schrödinger pourrait incorporer,
induire la réduction du vecteur d’état (voir par exemple Pearle, ou encore Ghirardi, Rimini
et Weber (GRW), pour des modèles sans variables supplémentaires prises au sens de la
Section 4 précédente). Cette réduction pourrait se produire à des moments récurrents (entre
lesquels une superposition indéfinie physiquement subsisterait) mais aléatoires suite aux
conditions expérimentales microscopiques non reproductibles. Il pourrait donc exister un
terme additionnel dans l’Équation de Schrödinger de nature fluctuante, stochastique. Dans
ce contexte, l’observation ne jouerait alors plus de rôle particulier autre que de capter,
mettre en évidence des propriétés physiques déjà bien définies via l’Équation de
Schrödinger, à certains instants seulement. La présente approche peut donc être qualifiée de
réaliste, dans une version du réalisme avec vérité connaissable (c.f. 1er tableau). Plus
concrètement, par exemple dans le modèle de GRW, des « localisations spontanées » ont
lieu dans l’espace, des sauts brusques récurrents avec une précision imparfaite
 {\footnotesize  Mais inférieure à l’incertitude expérimentale de la nouvelle ‘mesure’ (devenue simple prise de donnée) afin de
pouvoir retrouver un résultat unique – comme dans toute expérience.} (certaine
largeur de la fonction d’onde piquée en un point) mais ne permettant pas de superposition au
niveau macroscopique (décohérence spatiale assurée du même coup). Les probabilités de
localisation obtenues dans ce modèle sont normalisées sur tout l’espace [plus précisément
sur les positions possibles rj indicées par l’entier j] avec une amplitude plus élevée dans les
régions où le système ponctuel a déjà une plus forte probabilité de présence, ce qui a pour
effet de le localiser à terme en un point – lui ajoutant successivement des résonances (pics)
sur sa fonction d’onde restant normalisée au cours du temps. L’implémentation de ce
modèle dans l’Équation de Schrödinger se réalise via le nouveau terme ci-dessous,
\[
i\hbar\frac{d\rho(t)}{dt} = \lambda\sum_j[F_j\rho(t)F_j-\rho(t)]+[H,\rho(t)]
\]
où t désigne le temps, $\lambda$ une probabilité par unité de temps, H l’Hamiltonien, F$_j$(r)
l’opérateur localisant (par exemple une courbe gaussienne en r centrée en r$_j$) la fonction
d’onde $\ph{r}{\Psi}$ et $\rho$(t) l’opérateur densité :
\[
\rho(t)=\ket{\Psi(t)}\bra{\Psi(t)}
\]
Cette équation correspond à l’Équation de Schrödinger (dans sa version impliquant
l’opérateur densité $\rho$(t) plutôt que l’état associé $\ket{\Psi(t)}$ ) à laquelle a été
 ajouté le terme qui
succède juste au symbole d’égalité. Ce terme représente une évolution temporelle possible
de $\rho$(t) vers $F_j\ket{\Psi(t)}\bra{\Psi(t)}F_j$ . Notons qu’il existe des extensions
 continues du modèle
GRW, au prix de nouveaux termes encore, et même relativistes (ainsi qu’en théorie
quantique des champs). De plus, la réalisation de la localisation spontanée de GRW pourrait
trouver son origine dans la gravité (Diosi) via une localisation stochastique de la masse –
dont l’intensité serait proportionnelle à la constante universelle de gravitation (G) de
I.Newton. Cette voie rejoint celle de R.Penrose qui écrivait : « le seuil quantique de la
gravité qui produit une auto-réduction du vecteur d’état joue un rôle vis-à-vis de la
conscience » – possibilité plus générique donc d’une réduction du paquet d’onde par des
effets de Relativité Générale (voir aussi Bricmont). Enfin, une possibilité expérimentale de
(in)valider le modèle GRW serait la détection du faible effet de chauffage [atomes excités
rayonnant] spontané prévu par les localisations spontanées mêlées à la dynamique
temporelle (amenant même des questions sur la conservation de l’énergie lorsque combinées
à des réductions d’états d’impulsions).
\subsection{Familles cohérentes (ajout d’opérateurs)}
Considérons un état initial au temps t$_0$ associé à l’opérateur densité
\[
\rho(t)=\ket{\Psi(t)}\bra{\Psi(t)}
\]
La probabilité P($\mr{E}_1\mr{E}_2\mr{E}_3$...) que cet état appartienne au sous-espace de
Hilbert $\mr{E}_1$ au temps t$_1$ puis au sous-espace $\mr{E}_2$ au temps t$_2$ etc peut
être obtenue par la combinaison d’amplitudes de
transitions et constitue finalement la formule de Wigner :
\[
P(\mr{E}_1\mr{E}_2\mr{E}_3...)=
Tr\{...P_3(t_3)P_2(t_2)P_1(t_1)\rho(t_0)P_1(t_1)P_2(t_2)P_3(t_3)...\}
\]
où $P_i(t_i)$ est le projecteur orthogonal [de type $\sum_{k=1}^D\ket{u_k}\bra{u_k}$ à un
 instant donné] dans $\mr{E}_i$ –
de dimension D – au temps t$_i$ (dépendance temporelle dans l’opérateur selon le point de vue
de Heisenberg). La somme des probabilités, plutôt que des amplitudes, peut être bien
normée,
\[
\sum_{\{\mr{E}\}}P(\mr{E}_1\mr{E}_2\mr{E}_3...)=1
\]
si l’on ne considère en effet que les « histoires cohérentes » c’est-à-dire les états consécutifs
n’étant pas des combinaisons linéaires d’états appartenant à des sous-espaces Ei différents.
L’ensemble des histoires cohérentes forme, par définition (Griffiths puis Omnès et Gell-
Mann), une famille cohérente (définie donc par un ensemble de temps et d’un ‘découpage’
spécifique plus ou moins précis de l’espace pour chacun de ces temps). Les états
successivement projetés constituant ces histoires, ayant la vertu de ne pas être des
superpositions, pourraient représenter les véritables états du système à considérer au cours
du temps et les projecteurs associés (hermitiens) des observables ‘physiques’ – aux valeurs
propres 1 pour chacun de ces états à un temps donné et 0 pour tout état de l’espace
complémentaire en question. À remarquer néanmoins que lors de la transition entre deux
espaces consécutifs, seule une partie du vecteur d’état projeté subsiste, mais sans que cela
n’apparaisse tout-à-fait explicitement et formellement dans l’expression de
$P(\mr{E}_1\mr{E}_2\mr{E}_3...)$.
Quant aux séries de mesures (aux amplitudes de probabilité génériques bien normalisées)
possibles ‘en parallèle’ des projections, entre des mêmes temps initial t$_0$ et final t$_k$ , celles-ci
sont dictées par le formalisme usuel de la MQ, de manière compatible avec les résultats
expérimentaux obtenus sur les positions, impulsions... représentant elles une catégorie
d’observables habituelles. Ces séries de mesures doivent avoir une certaine correspondance
avec les projections successives – correspondance notamment fixée par le choix de la classe
de famille cohérente adéquate. En conclusion, les véritables états quantiques d’un système
au cours d’une période donnée, période de réalisation générale de mesures successives
(période comprenant donc notamment les phases antérieures aux mesures), sont
systématiquement définis de manière univoque (sans superposition), ce qui justifie ici une
approche réaliste. Il s’agit plus précisément de la version du réalisme avec vérité
inconnaissable puisque les états ‘pré-mesure’ sont possiblement des combinaisons linéaires
d’états propres des observables mesurables (habituelles).
\subsection{Univers multiple (nouvelles applications des lois quantiques)}
Le postulat de réduction du paquet d’onde est évité si l’on (sup)pose que chacune des
valeurs, associées à tous les vecteurs d’état d’une combinaison linéaire, est effectivement
obtenue lors d’une mesure (Everett, Graham, Saunders, Wallace et approche adoptée par
DeWitt pour qui : « le formalisme mathématique de la théorie quantique est capable de
fournir sa propre interprétation »). La superposition quantique est alors véritablement
réalisée (!). Il faut alors faire l’une des deux hypothèses fortes suivantes.

(i) Soit l’observateur possède lui-même plusieurs ‘composantes’ (comme une conscience
et plusieurs inconscients), chacune révélant l’un des résultats quantitatifs mais ne pouvant le
communiquer aux autres composantes – donnant bien l’impression d’un résultat unique dans
toute expérience. Et la probabilité d’obtenir un certain résultat est ici interprétée comme la
probabilité de se situer dans une certaine composante. Ainsi, la répétition d’une même
mesure d’observable sur un système préparé de manière identique génère toujours plusieurs
résultats selon la loi de probabilité de Born [déterminisme pour une composante donnée
mais non pour l’observateur], mais des résultats variés s’interprétant maintenant via des
changements de composantes pour l’observateur. Dans une composante, l’acte de mesure ne
fait que révéler un résultat (valeur propre) associé à un vecteur (propre) d’état pré-existant
vis-à-vis de la procédure de mesure – puisque tous les vecteurs propres revêtent ici une
réalité après et donc également avant la mesure (par continuité usuelle des états). Par
conséquent, on peut se trouver dans une vision réaliste avec vérité indépendante accessible.
Il s’agit-là de l’interprétation dite des « multivers », dans un monde restant unique.

(ii) Soit ces diverses composantes sont en fait des univers différents et la description ci-
dessus est transposée en conséquence. Il s’agit alors de l’« interprétation des mondes
multiples ». Ce scénario représente une réalisation de mondes parallèles permettant de
contourner partiellement le « principe anthropique fort » (= réglage des constantes
fondamentales de la nature dans l’objectif d’une apparition de la vie) : chaque univers
parallèle peut correspondre à une valeur possible (associée à un état d’une combinaison
linéaire dont le nombre de termes peut augmenter avec les assemblages d’atomes et autres
interactions implémentées dans l’équation de Schrödinger) des nombreuses observables (les
résultats des mesures humaines correspondant à des sous-classes de l’ensemble complet
d’univers) et la probabilité est grande que dans au moins quelques univers la vie apparaisse
bien. Nous sommes par définition dans l’un de ceux-ci. Néanmoins ici les constantes
fondamentales elles permettent l’émergence d’organismes vivants dans au moins un
univers ; un contournement plus radical du principe anthropique fort serait que les valeurs
de ces constantes fondamentales soient distribuées au hasard parmi un grand nombre
d’univers parallèles de telle sorte que la probabilité d’apparition de la vie dans l’un d’eux
soit élevée. Notons enfin que dans le contexte de la recherche en cosmologie, surgit aussi
parfois [questions de la matière/énergie noire, des transitions de phase, de l’inflation,...]
l’idée de la présence de plusieurs univers (aux conditions physiques différentes) mais
correspondant plutôt à différents domaines d’un monde unique.

Les deux interprétations ci-dessus doivent, d’une manière ou d’une autre, avoir un caractère
collectif au vu de l’universalité des expériences/expérimentations scientifiques
reproductibles. Cette universalité est par exemple explicitement respectée dans la version de
multivers appelée solipsisme [philosophie selon laquelle la seule réalité est le sujet pensant
lui-même, donc lié à ses observations] convivial (Zwirn).

